\subsubsection{二分图最小权完美匹配 (KM / SSP)}
Kuhn-Munkres (KM) 算法，用于 O(n\textasciicircum 3) 求解二分图最小/大权完美匹配。
\begin{lstlisting}
ll cst[1010][1010];
#define inf 0x3f3f3f3f3f3f3f3f
pair<ll,vector<int>> KM(int n,int m){
    if(n>m)return {inf,{}};
    vector<ll>l(n+1),r(m+1);
    vector<int>p(m+1),ans(n+1);
    f(i,1,n){
        vector<ll>dis(m+1,inf);
        vector<int>pre(m+1),done(m+1);
        int x,v=0,u=0;
        for(p[0]=i;x=p[u];u=v){
            done[u]=1;
            ll mn=inf;
            f(j,1,m)if(!done[j]){
                ll w=cst[x][j]-l[x]-r[j];
                if(w<dis[j])dis[j]=w,pre[j]=u;
                if(dis[j]<mn)mn=dis[j],v=j;
            }
            if(mn==inf)return {inf,{}};
            f(j,0,m){
                if(done[j])l[p[j]]+=mn,r[j]-=mn;
                else dis[j]-=mn;
            }
        }
        for(int v;u;u=v)v=pre[u],p[u]=p[v];
    }f(j,1,m)if(p[j])ans[p[j]]=j;
    return {-r[0],ans};
    // cout<<(-r[0])<<endl;//最小权
    // f(j,1,m)cout<<p[j]<<" ";cout<<endl;//右部点的匹配点。左部点匹配点是 ans。
}
\end{lstlisting}

\paragraph{稀疏图费用流做法（SSP）}
复用前文 MCMF 代码而不是另写一份松弛器。下例为与其放在同一作用域的适配函数，左右点分别从 1 编号；每组重新建残量网络。最大匹配数下求最小费用，要求完美匹配时额外检查返回流量是否等于左部点数。允许原始边负费用；初始网络为 DAG，因此无负费用环。
\begin{lstlisting}
array<ll,2> sparse_matching(int left,int right,
                          const vector<tuple<int,int,ll>>& edges){
    nodes=left+right+2;S=left+right+1;T=S+1;
    assert(nodes<100010&&2*(left+right+edges.size())+2<1000010);
    fill(hd,hd+nodes+1,0);cnt=1;cst=0;
    for(int u=1;u<=left;u++)ad(S,u,1,0);
    for(int v=1;v<=right;v++)ad(left+v,T,1,0);
    for(auto [u,v,c]:edges)ad(u,left+v,1,c);
    spfa(S);return mcmf();
}
\end{lstlisting}

\paragraph{使用条件与多测} 稠密 KM 与稀疏 SSP 为两个接口。第一份返回最小费用及左部匹配，要求 $n\le m$、完整有限费用矩阵；每组必须完整覆盖 cst[1..n][1..m]，其余状态局部化。不可用边不要直接当无限大参与减势运算，稀疏图优先用费用流；若用有限哨兵需证明足够大且最终检查是否选到禁边。矩形复杂度 $O(n^2m)$。SSP 适配器依赖前文 MCMF（去掉其 main 后合用），会重建流网络；费用和减势运算不得溢出，目标是最大匹配数下的最小费用，未必完美匹配。复杂度为初始化加每次增广的 $O(E\log V)$。
